\scanpage{58}

% manuscript line(s) 1
Donc $X'/X$ est un \uline{torseur} sous $G_T$ ! Et \unsure{ce|se|si|2}

% manuscript line(s) 2–6
\uline{Th} On a une équivalence de catégories \inserted{(groupoïdes)} fibrées sur $\mathrm{Sch}/\mathbb F_2$ (plus gén. sur les topos loc\textsuperscript{t} annelés en car. $2$), \unsure{entre} la catégorie des rev. quadratiques relatifs (sur $X\in\mathrm{Ob}\ \mathrm{Sch}/\mathbb F_2$), et les données suivantes sur un tel $X$
% manuscript line(s) 7–10
\[
(L,\inserted{P_1,}T,\alpha),\quad\text{où}\quad\left\{\begin{aligned}&L\text{ module inv. sur }X,\ P_1\ L\text{-tors.}\\ &T\text{ section de }L\ \text{(module des translations de }P_1\text{)}\\ &\alpha\text{ section de}\\ &\quad P_1\wedge_{\underline L}(L^{\otimes2},f_T:L\unsure{\mathbin{-}|\to}L^{\otimes2})\end{aligned}\right.
\]
% manuscript line(s) 7
\margin{(i.e. $P_1$ un fibré en droites \uline{affines} sur $X$)}
% manuscript line(s) 11
où $f_T:L\mathbin{-}L^{\otimes2}$ donné par $f_T(u)=u^2+T\cdot u$
% manuscript line(s) 12–13
(c'est un \unsure{homo|hom.} additif $W(L)\mathbin{-}W(L^{\otimes2})$, dont le noyau est un \inserted{schéma en} \unsure{groupe|groupes} \inserted{fini} quadratique $G_T$ sur $X$).
% manuscript line(s) 14
Au système $\Sigma=(L,P_1,T,\alpha)$ est associé le $X$-schéma
% manuscript line(s) 15
\[
X'_\Sigma=f_{T,P_1}\unsure{\zeta^{-1}|{}^{-1}}(V(\alpha))
\]
% manuscript line(s) 16
(torseur sous $G_T$), qui est \inserted{fini} quadratique sur $X$.

% manuscript line(s) 17
Plus de doute, c'est lui !
% manuscript line(s) 17–22
\margin{\uline{NB} $u:P_1\longrightarrow P_1^{(2)}$ comp. avec $u\mapsto Tu$ est donné par $u:x\longmapsto Tx+N$, ou $\underline{u(x)=x^2+Q(x)}$ en effet $(Q(x)=x^2+Tx+N)$, donc \deleted{\unsure{on aura}} \deleted{\unsure{$X'$}} l'équation $Q(x)=0$ qui car. $X'/X$ équivaut à $\boxed{u(x)=x^2}$ \quad $P_1\mathrel{\substack{\xrightarrow{u}\\[-3pt]\xrightarrow[\ x\mapsto x^2\ ]{\phantom{u}}}}P_1^{(2)}$ \deleted{\unsure{$f$}}}

% manuscript line(s) 18–19
Comment s'explicite le produit $\wedge$ des rev. quadratiques ?
% manuscript line(s) 20–21
\[
\begin{aligned}(L,T,P_1,\alpha)\wedge(L',T',P'_1,\alpha')&\overset?={}\\ &(L\otimes L',TT',?,?)\end{aligned}
\]
% manuscript line(s) 22–23
\unsure{Ainsi|Alors} :
% manuscript line(s) 22–23
\[
\begin{aligned}G_T\xrightarrow{f_{T'}}G_{TT'}&\quad\text{induit par } f_{T'}:L\mathbin{-}L\otimes L',\quad u\longmapsto uT'\\G_{T'}\mathbin{-}G_{TT'}&\quad\rule{1cm}{.3pt}\ f_T:L'\mathbin{-}L\otimes L',\quad u'\mathbin{-}Tu'\end{aligned}
\]
% manuscript line(s) 24–26
d'où $G_T\times G_{T'}\longrightarrow G_{TT'}$ hom de groupes d'où, pour des torseurs $\unsure{\mathscr X|\mathcal X|\mathfrak X}$ sous $G_T$, $\mathscr X'$ sous $G_{T'}$, un torseur $\mathscr X\wedge\mathscr X'$ sous $G_{TT'}$ !
% archive sheet number: 57
